\chapter{Software and Resource Availability}\label{app:software}

To support open scientific reproducibility, transparent peer review, and community adoption across academic proteomics laboratories and biotechnology pipelines, all computational resources developed in this work---including the neural network source code, dynamic programming algorithms, pre-trained model weights, evaluation pipelines, and interactive deployment platforms---are publicly available under permissive open-source licenses.

\section{Open-Source Codebase and Repository Architecture}

The official implementation of \textbf{DFlowNovo} is hosted on GitHub under the permissive Apache~2.0 license:
\begin{center}
\url{https://github.com/joelinator/JoelResearch}
\end{center}

The repository is organized modularly to enable standalone usage of individual algorithmic components:
\begin{itemize}
    \item \texttt{src/models/}: PyTorch implementations of the bidirectional spectrum Transformer encoder, auxiliary peptide length predictor, and discrete flow matching decoder with AdaLN-Zero conditioning.
    \item \texttt{src/algorithms/}: The exact dynamic programming knapsack reachability engine (\texttt{KnapsackDP}), precomputing discrete residue prefix and suffix validity tables and executing $\mathcal{O}(1)$ GPU bitmask lookups.
    \item \texttt{src/data/}: Modular spectrum parsers and loaders supporting standard mass spectrometry formats, including Mascot Generic Format (\texttt{.mgf}) and compressed Apache Parquet tables.
    \item \texttt{scripts/infer.py}: Command-line interface for batch peptide sequencing from unannotated experimental MS/MS spectra.
    \item \texttt{scripts/eval.py}: Standardized evaluation suite computing strict exact match, isobaric Leucine/Isoleucine match, residue-level precision/recall/F1, and identification coverage curves.
    \item \texttt{scripts/generate\_publication\_figures.py}: Automated plotting routines generating all high-resolution figures in this thesis.
    \item \texttt{tests/}: Comprehensive unit test suite covering spectrum featurization, rate matrix dynamics, dynamic programming boundary conditions, and end-to-end model inference.
\end{itemize}

\subsection{Installation and Environment Reproduction}

The environment requires Python 3.10 or higher and PyTorch 2.1+ with CUDA support:
\begin{verbatim}
git clone https://github.com/joelinator/JoelResearch.git
cd JoelResearch
python3 -m venv .venv
source .venv/bin/activate
pip install -r requirements.txt
pytest tests/
\end{verbatim}

\section{Pre-Trained Model Weights and Checkpoints}

The production model weights for DFlowNovo are hosted on the Hugging Face Model Hub:
\begin{center}
\url{https://huggingface.co/joelinator/dflow-novo-model}
\end{center}

\subsection{Checkpoint Specifications}
\begin{itemize}
    \item \textbf{File Name:} \texttt{frozen\_production\_model.ckpt}
    \item \textbf{File Size:} 453.9\,MB (optimized for fast distribution and containerization)
    \item \textbf{Total Parameters:} 59,476,250
    \item \textbf{Numerical Precision:} Bfloat16 (BF16) weights with FP32 KnapsackDP integer masks
    \item \textbf{SHA-256 Checksum:} \texttt{4fa8639ce8d56b0632b6db760ca4a9b69fc03f90bcfc70f5e55e8211db4850fa}
\end{itemize}

The checkpoint can be downloaded automatically via the Hugging Face Hub CLI:
\begin{verbatim}
huggingface-cli download joelinator/dflow-novo-model \
    frozen_production_model.ckpt --local-dir models/
\end{verbatim}

\section{Interactive Cloud Demonstration Studio}

An interactive, zero-installation web demonstration powered by Gradio is hosted on Hugging Face Spaces:
\begin{center}
\url{https://huggingface.co/spaces/joelinator/dflow-novo}
\end{center}

The web interface enables users to:
\begin{enumerate}
    \item Upload custom experimental tandem mass spectrometry files (\texttt{.mgf}).
    \item Inspect experimental spectrum peaks, precursor neutral mass, and precursor charge state.
    \item Execute real-time DFlowNovo discrete flow decoding with adjustable Euler integration step budgets ($K \in [5, 30]$).
    \item Toggle dynamic programming KnapsackDP mass guidance to observe physical filtering effects.
    \item Visualize theoretical fragment ion mirror plots ($b$- and $y$-ions) aligned against experimental peaks.
\end{enumerate}

\section{Standardized Benchmark Proteomics Datasets}

All training, validation, and test splits used throughout this thesis were ingested from public datasets released by InstaDeep on the Hugging Face Hub:
\begin{itemize}
    \item \textbf{Nine-Species Biological Benchmark:}
    \begin{center}
    \url{https://huggingface.co/datasets/InstaDeepAI/ms_ninespecies_benchmark}
    \end{center}
    Originally described by \citet{Tran2017}, the raw experimental mass spectrometry files are preserved at the UCSD MassIVE repository (Accession: \texttt{MSV000081382}) and the ProteomeXchange Consortium (Accession: \texttt{PXD004948}).
    \item \textbf{ProteomeTools HC-PT Synthetic Benchmark:}
    \begin{center}
    \url{https://huggingface.co/datasets/InstaDeepAI/ms_proteometools}
    \end{center}
    Originally generated by \citet{Zolg2017}, the raw acquisitions reside in the EMBL-EBI PRIDE database (Accession: \texttt{PXD004732}) and UCSD MassIVE (Accession: \texttt{MSV000080502}).
\end{itemize}

\section{Computational Infrastructure and Hardware Environment}

Model training and evaluation benchmarks reported in Chapter~\ref{chap:benchmarks} were executed using the following hardware and software specifications:
\begin{itemize}
    \item \textbf{Accelerators:} $4\times$ NVIDIA H100 SXM5 GPUs (80\,GB HBM3 memory per GPU) for distributed multi-GPU training; single NVIDIA RTX 4090 / H100 for inference throughput benchmarks.
    \item \textbf{Host Compute:} 64-core AMD EPYC 7763 CPU @ 2.45\,GHz, 512\,GB DDR4 RAM.
    \item \textbf{Software Stack:} Ubuntu 22.04 LTS, CUDA 12.2, cuDNN 8.9, PyTorch 2.1.2, FlashAttention-2.
\end{itemize}
